\section{New 2025 courses: Geometric DL, Unsupervised, DL Architectures}

The 2025 curriculum added \emph{Deep Learning: Geometric Techniques} and renamed the deep learning
courses (\emph{Basic Techniques} = old DLNN~I, \emph{Architectures and Generative Techniques} =
old DLNN~II). An examiner who teaches them may use their vocabulary, so this chapter covers what is
new. Most of DL Basic is already in Chapters~1--3.

\subsection{Deep Learning: Geometric Techniques}

\begin{summary}{Geometric deep learning in two minutes}
\textsf{Problem:} the curse of dimensionality; local methods need exponentially many samples in
high dimension. \textsf{Remedy:} priors about the \emph{structure} of the data. \
\textsf{Two priors.} (1) \emph{Symmetry}: the task does not change under a group $G$ of
transformations (translations for images, permutations for sets and graphs, rotations for
molecules). (2) \emph{Scale separation}: coarse-grain the signal step by step. \
\textsf{Invariance} $f(\rho(g)x)=f(x)$, \textsf{equivariance} $f(\rho(g)x)=\rho'(g)f(x)$. \
\textsf{Blueprint:} equivariant layers $\to$ non-linearity $\to$ local pooling (coarsening) $\to$
repeat $\to$ global invariant pooling $\to$ MLP. CNNs (grids, translations), DeepSets (sets,
permutations), GNNs (graphs, permutations), Transformers (complete graphs) are all instances. \
\textsf{Message passing:} $\mathbf{h}_i'=\phi\bigl(\mathbf{h}_i,\bigoplus_{j\in\mathcal{N}(i)}\psi(\mathbf{h}_i,\mathbf{h}_j,\mathbf{e}_{ij})\bigr)$,
$\bigoplus$ permutation-invariant (sum, mean, max).
\end{summary}

\pic{A face is still a face if you move the photo (translation) and a molecule is still the same
molecule if you number its atoms differently (permutation). Building these facts into the network
means it never has to learn them from data, which is the cure for the curse of dimensionality.}

\board{Group axioms: closure, associativity, identity, inverse \textbullet{} equivariance
$f(\rho(g)x)=\rho'(g)f(x)$ \textbullet{} DeepSets $f(X)=\rho\bigl(\sum_i\phi(\mathbf{x}_i)\bigr)$ \textbullet{}
set permutation $X\mapsto PX$ \textbullet{} GCN $H'=\sigma(\hat D^{-1/2}\hat A\hat D^{-1/2}HW)$ \textbullet{}
group averaging $\bar f(x)=\frac{1}{|G|}\sum_{g\in G}f(\rho(g)x)$.}

\begin{qabox}[GDL lecturer, Klambauer]{\cexam What is the curse of dimensionality, and what is the manifold hypothesis?}
\from{DL: Geometric Techniques · curse of dimensionality · course exam question map}
\ans In high dimension the volume grows so fast that data become sparse; local methods (nearest
neighbours, kernels with small bandwidth) need a number of samples exponential in the dimension.
For Lipschitz functions the estimation error decays only like $N^{-1/d}$. The manifold hypothesis
says real data lie near a low-dimensional manifold inside the high-dimensional space, so the
\emph{intrinsic} dimension is small; together with symmetry priors this makes learning possible.
\end{qabox}

\begin{qabox}[GDL lecturer]{\cexam Define a group, invariance and equivariance.}
\from{DL: Geometric Techniques · geometric priors · course exam question map}
\ans A group is a set with an operation that is closed and associative, has an identity, and has an
inverse for every element. It acts on the domain, and the action lifts to signals via a
representation $\rho(g)$ (a linear map). $f$ is invariant if $f(\rho(g)x)=f(x)$, equivariant if
$f(\rho(g)x)=\rho'(g)f(x)$: transform then map equals map then transform (a commutative diagram).
Inverses are needed so that the lifted action $(\rho(g)x)(u)=x(g^{-1}u)$ is well defined.
\follow ``How do you make any function invariant?'' Average over the group:
$\bar f(x)=\frac1{|G|}\sum_gf(\rho(g)x)$. The approximation error does not grow if the target is
invariant; the estimation error shrinks, roughly by the size of the group.
\end{qabox}

\begin{qabox}[GDL lecturer, Klambauer]{\cexam DeepSets and linear equivariant maps on sets.}
\from{DL: Geometric Techniques · sets · course exam question map}
\ans The symmetry group of a set of $n$ elements is the permutation group $S_n$; a permutation acts
on the feature matrix as $X\mapsto PX$. Every linear permutation-equivariant layer has the form
$F(X)=XA+\frac1n\mathbf{1}\mathbf{1}^\top XB$ (the same map for every element plus a shared term).
DeepSets: an equivariant per-element encoder $\phi$ (learned), a sum (fixed, invariant), then a
learned decoder $\rho$: $f(X)=\rho(\sum_i\phi(\mathbf{x}_i))$.
\end{qabox}

\begin{qabox}[GDL lecturer, Klambauer]{\cexam Message passing, and the Weisfeiler--Leman test.}
\from{DL: Geometric Techniques · graphs; also asked (GNNs) in past defenses}
\ans A graph adds edges to a set. Each layer: every node computes messages from its neighbours,
aggregates them with a permutation-invariant function, and updates its state. Node outputs are
equivariant, a final readout makes graph outputs invariant. The WL test refines node colours by
hashing each node's colour with the multiset of its neighbours' colours; if the histograms differ,
the graphs are not isomorphic; if they agree, we cannot tell. MPNNs are at most as powerful as WL,
and exactly as powerful if the aggregation is injective on multisets (sum, as in GIN, not mean or max).
\follow ``Limits?'' Over-squashing: information from many distant nodes is squeezed through few
edges; it depends mostly on graph topology (bottlenecks). Remedies: positional encodings,
subgraph GNNs, graph rewiring.
\end{qabox}

\begin{qabox}[GDL lecturer]{\cexam Transformers as graph networks; CNNs as the grid case.}
\from{DL: Geometric Techniques · graphs, grids · course exam question map}
\ans Self-attention is message passing on the complete graph with attention weights as edge
weights. Without positional encodings it is permutation-equivariant; with them it is not, by
design. A CNN is the blueprint on a grid with the translation group: convolution is the (only)
linear translation-equivariant map, pooling does the coarse-graining.
\end{qabox}

\subsection{ML: Unsupervised Techniques}

\begin{summary}{Unsupervised learning in two minutes}
\textsf{PCA:} directions of maximal variance = eigenvectors of the covariance; unique up to sign if
eigenvalues differ. \textsf{Kernel PCA:} PCA on the centred Gram matrix
$\tilde K=K-\mathbf{1}_NK-K\mathbf{1}_N+\mathbf{1}_NK\mathbf{1}_N$. \
\textsf{ICA:} find \emph{statistically independent} non-Gaussian sources; unique up to order and
scale; whitening then rotation; contrast functions (kurtosis, negentropy); Infomax. \
\textsf{Factor analysis:} $\mathbf{x}=W\mathbf{z}+\bm\mu+\bm\epsilon$ with diagonal (not isotropic)
noise, so $\mathbf{x}\sim\mathcal{N}(\bm\mu,WW^\top+\Psi)$; fitted with EM. \
\textsf{MDS:} embed points so distances are kept. \textsf{NMF:} $X\approx WH$, $W,H\ge0$, parts-based. \
\textsf{Mixtures + EM}, \textsf{clustering} (k-means = hard EM), \textsf{max entropy}.
\end{summary}

\pic{PCA finds the longest axis of the cloud of points. ICA is the cocktail party: two people talk
at once into two microphones, and ICA separates the two voices because each voice on its own is
not Gaussian.}

\board{PCA: $C=\frac1N\sum_n(\mathbf{x}_n-\bar{\mathbf{x}})(\mathbf{x}_n-\bar{\mathbf{x}})^\top$, $C\mathbf{u}=\lambda\mathbf{u}$ \textbullet{}
FA: $\mathbf{x}\sim\mathcal{N}(\bm\mu,WW^\top+\Psi)$ \textbullet{}
ICA: $\mathbf{x}=A\mathbf{s}$, $\mathbf{y}=W\mathbf{x}$, maximize non-Gaussianity \textbullet{}
k-means: $\min\sum_n\lVert\mathbf{x}_n-\bm\mu_{c(n)}\rVert^2$.}

\begin{qabox}[Klambauer, Hochreiter]{\cexam What is PCA, how is it computed, is it unique?}
\from{ML: Unsupervised Techniques · PCA · Fragenkatalog Q8--10}
\ans Centre the data, compute the covariance matrix, and take its eigenvectors with the largest
eigenvalues (or the SVD of the data matrix). They are the orthogonal directions of maximal
variance, and projecting on them gives the best linear reconstruction in squared error. Unique up
to the sign of each vector when the eigenvalues are distinct; with equal eigenvalues any rotation
inside that eigenspace is also a solution. Properties: linear, orthogonal, uncorrelated components.
\end{qabox}

\begin{qabox}[Klambauer, Hochreiter]{\cexam ICA versus PCA versus factor analysis.}
\from{ML: Unsupervised Techniques · ICA, FA · Fragenkatalog Q13--23}
\ans PCA: uncorrelated, orthogonal directions of maximal variance. ICA: statistically
\emph{independent} components, which needs non-Gaussian sources; unique up to permutation and
scaling. It whitens the data (then only a rotation is left) and finds the rotation that maximizes
non-Gaussianity with a contrast function such as kurtosis or negentropy; Infomax maximizes the
information flow through a non-linearity. Factor analysis: a latent Gaussian model like pPCA but
with a diagonal noise covariance, so every feature has its own noise; fitted with EM, whose E-step
estimates the posterior mean and covariance of the factors.
\end{qabox}

\begin{qabox}[Klambauer]{\pred Clustering: k-means and its relation to EM.}
\from{ML: Unsupervised Techniques · clustering, mixtures · predicted}
\ans k-means alternates two steps: assign each point to the nearest centre, then move each centre
to the mean of its points; the objective $\sum_n\lVert\mathbf{x}_n-\bm\mu_{c(n)}\rVert^2$ never
increases. It is the hard-assignment limit of EM for a Gaussian mixture with equal isotropic
variances: responsibilities become 0 or 1.
\end{qabox}

\subsection{Deep Learning: Architectures and Generative Techniques}

\begin{summary}{DL Architectures in two minutes}
\textsf{CNN families:} LeNet $\to$ AlexNet $\to$ VGG (small $3\times3$ filters) $\to$ Inception
(parallel filters) $\to$ ResNet (skips) $\to$ DenseNet. \
\textsf{Vision tasks:} classification; semantic segmentation (class per pixel, U-Net, FCN);
object detection (boxes: R-CNN family two-stage, YOLO/SSD one-stage, IoU, non-maximum suppression);
instance segmentation (Mask R-CNN). \
\textsf{Generative:} autoencoder, VAE (ELBO), GAN (minimax), normalizing flows (exact likelihood by
change of variables), diffusion, autoregressive models (PixelCNN). \
$\min_G\max_D\ \mathbb{E}[\log D(\mathbf{x})]+\mathbb{E}[\log(1-D(G(\mathbf{z})))]$;\quad
flow: $\log p(\mathbf{x})=\log p(\mathbf{z})+\log\lvert\det\partial\mathbf{z}/\partial\mathbf{x}\rvert$.
\end{summary}

\pic{GAN: a forger (generator) and a detective (discriminator) train against each other until the
detective can no longer tell the fakes from the real paintings.}

\begin{qabox}[Klambauer]{\pred Semantic segmentation versus object detection. Name architectures.}
\from{DL: Architectures (DLNN II) · segmentation and detection, SS25 script PDF p.~197--239; U-Net asked in a past defense}
\ans Semantic segmentation gives every pixel a class; U-Net (encoder--decoder with skip
connections) or FCN. Object detection gives boxes plus classes; two-stage detectors (Faster R-CNN:
region proposals, then classification) or one-stage (YOLO, SSD: predict boxes on a grid directly).
Quality by intersection over union (IoU); duplicate boxes removed by non-maximum suppression.
Instance segmentation (Mask R-CNN) separates individual objects of the same class.
\end{qabox}

\begin{qabox}[Klambauer]{\pred GANs: objective, problems, fixes.}
\from{DL: Architectures (DLNN II) · GANs, SS25 script PDF p.~265--279 · predicted}
\ans The generator maps noise to samples, the discriminator tells real from fake; they play
$\min_G\max_D\mathbb{E}[\log D(\mathbf{x})]+\mathbb{E}[\log(1-D(G(\mathbf{z})))]$. At the optimum the
generator minimizes the Jensen--Shannon divergence. Problems: mode collapse, unstable training,
vanishing generator gradients. Fixes: non-saturating loss, Wasserstein GAN, spectral normalization,
and the two time-scale update rule (TTUR, from Linz), with the FID to measure quality.
\end{qabox}

\begin{qabox}[Klambauer]{\pred Autoencoder, VAE, normalizing flow: what is the difference?}
\from{DL: Architectures (DLNN II) · generative models, SS25 script PDF p.~242--297 · predicted}
\ans An autoencoder compresses and reconstructs; its latent space is not organised, so it is poor
for sampling. A VAE makes the latent space a distribution with a KL term, so we can sample, but the
likelihood is only bounded (ELBO). A normalizing flow uses invertible layers and the change of
variables formula, so the likelihood is exact, at the price of equal dimensions and constrained
layers. Diffusion trades speed for quality with many denoising steps.
\end{qabox}
